\documentclass[11pt]{article}
\usepackage{mathblatt}
\usepackage{multicol}


\begin{document}
\pfheftstil
\pfheftkopf{Flaechen \textperiodcentered{} Lösungen}{}
\begin{pfloesung}{}
\lz{1.}{3a; a²; 2(a + b)}{}{}
\end{pfloesung}
\begin{pfloesung}{}
\lz{2.}{A = a · b, u = 2 · (a + b); A = g · h : 2; A = (a + c) : 2 · h; A = e · f : 2; A = $\pi$ · r²}{}{}
\end{pfloesung}
\begin{pfloesung}{}
\lz{3.}{r = 2,5 m; (20 $-$ 12) : 2 = 4 m}{}{}
\end{pfloesung}
\begin{pfloesung}{}
\lz{4.}{Umfang}{}{}
\end{pfloesung}
\begin{pfloesung}{}
\lz{5.}{Strecke von C senkrecht auf AB}{}{}
\end{pfloesung}
\begin{pfloesung}{}
\lz{6.}{A = (a + b) · c}{}{}
\end{pfloesung}
\begin{pfloesung}{}
\lz{7.}{A = d · e : 2}{Katheten = Grundseite und Höhe}{}
\end{pfloesung}
\begin{pfloesung}{}
\lz{8.}{u = 3 · a}{drei gleich lange Seiten}{}
\end{pfloesung}
\begin{pfloesung}{}
\lz{9.}{u = 16a}{}{}
\end{pfloesung}
\begin{pfloesung}{}
\lz{10.}{A = a · (b + c)}{Gesamthöhe b + c}{}
\end{pfloesung}
\begin{pfloesung}{}
\lz{11.}{Strecke von der oberen Seite senkrecht hinunter auf $g$, rechter Winkel am Fußpunkt}{}{}
\end{pfloesung}
\begin{pfloesung}{}
\lz{12.}{$\tfrac{a + c}{2} \cdot h$}{$A$}{}
\end{pfloesung}
\begin{pfloesung}{}
\lz{13.}{2000 cm²}{A = e · f : 2}{}
\end{pfloesung}
\begin{pfloesung}{}
\lz{14.}{10 cm}{u = 2 · (a + b)}{}
\end{pfloesung}
\begin{pfloesung}{}
\lz{15.}{900$\pi$ $\approx$ 2827,4 cm²}{G = $\pi$ · r² = $\pi$ · 30² = 900$\pi$ cm²; r² = 900 cm²}{}
\end{pfloesung}
\begin{pfloesung}{}
\lz{16.}{gegeben AD = 4,1 cm, $\delta$ = 21°, rechter Winkel bei E; AE mit Sinus, DE mit Kosinus, A = ½ · AE · DE}{AE = 4,1 · sin 21° $\approx$ 1,47 cm; DE = 4,1 · cos 21° $\approx$ 3,83 cm}{}
\end{pfloesung}
\begin{pfloesung}{}
\lz{17.}{1,1449$\pi$ $\approx$ 3,60 m²}{r = 1,07 m}{}
\end{pfloesung}
\begin{pfloesung}{}
\lz{18.}{h\_s = √(0,68² + 0,40²) = √0,6224 $\approx$ 0,789 m $\Rightarrow$ A = 0,4 · √0,6224 $\approx$ 0,32 m²}{h\_s = √0,6224 $\approx$ 0,789 m; A = 0,4 · √0,6224 $\approx$ 0,316 m²}{}
\end{pfloesung}
\begin{pfloesung}{}
\lz{19.}{0,975 m²; V = 0,975 · 1,2 = 1,17 m³ $\approx$ 1,2 m³; $\approx$ 320 kg}{A = ½ · 1,5 · 1,3; 10 \% von 1,2 m³ = 0,12 m³ $\mathrel{\widehat{=}}$ 120 kg}{}
\end{pfloesung}
\begin{pfloesung}{}
\lz{20.}{$\pi \cdot r^2$}{$A$}{}
\end{pfloesung}
\begin{pfloesung}{}
\lz{21.}{$4{,}5$ cm}{Durchmesser; $r$}{}
\end{pfloesung}
\begin{pfloesung}{}
\lz{22.}{6 cm}{}{}
\end{pfloesung}
\begin{pfloesung}{}
\lz{23.}{105 mm}{A = a · b $\Rightarrow$ b = 42 000 mm² : 400 mm}{}
\end{pfloesung}
\begin{pfloesung}{}
\lz{24.}{b = 5 cm}{}{}
\end{pfloesung}
\begin{pfloesung}{}
\lz{25.}{QC = 15 m}{½ · QC · 12 = 90}{}
\end{pfloesung}
\begin{pfloesung}{}
\lz{26.}{nein}{Plattenlänge 1,5 m < 1,65 m; 1,8 m² : 1,2 m = 1,5 m}{}
\end{pfloesung}
\begin{pfloesung}{}
\lz{27.}{h = 2 · 5225 : 190 = 55 cm > 50 cm $\Rightarrow$ Tiefe reicht}{(110 + 80) : 2 · h = 5225 $\Rightarrow$ h = 55 cm}{}
\end{pfloesung}
\begin{pfloesung}{}
\lz{28.}{ein Rechteck und ein rechtwinkliges Dreieck}{}{}
\end{pfloesung}
\begin{pfloesung}{}
\lz{29.}{$\tfrac{1}{2}$}{}{}
\end{pfloesung}
\begin{pfloesung}{}
\lz{30.}{Rechteck 10 m × 5 m und zwei Halbkreise (zusammen ein Kreis, r = 2,5 m)}{}{}
\end{pfloesung}
\begin{pfloesung}{}
\lz{31.}{$\tfrac{29}{72}$ $\approx$ 40,3 \%}{145 : 360 = $\tfrac{29}{72}$ $\approx$ 0,4028}{}
\end{pfloesung}
\begin{pfloesung}{}
\lz{32.}{A = 13 · 5 cm² + $\pi$ · 2,5² cm² $\approx$ 65 cm² + 19,63 cm² $\approx$ 84,63 cm²}{}{}
\end{pfloesung}
\begin{pfloesung}{}
\lz{33.}{240 $-$ 10,24$\pi$ $\approx$ 207,83 m²}{Rechteck 240 m²; Kreis $\pi$ · 3,2² = 10,24$\pi$ $\approx$ 32,17 m²}{}
\end{pfloesung}
\begin{pfloesung}{}
\lz{34.}{50 + 6,25$\pi$ $\approx$ 69,63 m²}{Rechteck 50 m²; Kreis 6,25$\pi$ $\approx$ 19,63 m²}{}
\end{pfloesung}
\begin{pfloesung}{}
\lz{35.}{A = 3,7 · (√144,05 + 7,4 : tan 52°) $\approx$ 65,8 m²}{DB = 7,4 : tan 52° $\approx$ 5,8 m; c = 12,0 + 5,8 $\approx$ 17,8 m}{}
\end{pfloesung}
\begin{pfloesung}{}
\lz{36.}{A\_Kreis $\approx$ 10 386,89 cm²; Verschnitt $\approx$ 4 613,11 cm² $\approx$ 30,75 \%}{}{}
\end{pfloesung}
\begin{pfloesung}{}
\lz{37.}{Ja}{12 Deckel, Abfall 800 $-$ 192$\pi$ $\approx$ 196,8 cm² $\approx$ 24,6 \% $\le$ 25 \%; 100 : 8 = 12,5 $\Rightarrow$ 12 Deckel; Deckelfläche 12 · $\pi$ · 16 = 192$\pi$ $\approx$ 603,2 cm²; Blech 800 cm²}{}
\end{pfloesung}
\begin{pfloesung}{\textbf{Prüfstein}}
\lz{a)}{$\gamma$ = 124°}{180° $-$ 56°}{}
\lz{b)}{A = 163,2 m²}{Mittellinie (25,8 + 15) : 2 = 20,4 m}{}
\lz{c)}{U = 40,8 + 16 : sin 56° $\approx$ 60,1 m}{Schenkel s = 8 : sin 56° $\approx$ 9,65 m (oder √(5,4² + 8²)); Überstand (25,8 $-$ 15) : 2 = 5,4 m}{}
\end{pfloesung}
\end{document}